% Источники: 05_Трассировка/DOC-09_Вклад_Самбьюи.md
% Генератор: docs/build_pdf.py
% Шаблон печатного представления. Содержимое подставляет docs/build_pdf.py.
\documentclass[12pt,a4paper]{article}
\usepackage{iftex}
\ifPDFTeX
  \usepackage[utf8]{inputenc}
  \usepackage{cmap}
  \usepackage[T2A]{fontenc}
  % Tempora retains Times-like text and provides Cyrillic bold/italic shapes.
  \usepackage{tempora,newtxmath}
\else
  \usepackage{fontspec}
  \IfFontExistsTF{Times New Roman}
    {\setmainfont{Times New Roman}}
    {\setmainfont{Libertinus Serif}}
  \IfFontExistsTF{Arial}{\setsansfont{Arial}}{\setsansfont{Latin Modern Sans}}
  \IfFontExistsTF{Menlo}{\setmonofont{Menlo}[Scale=MatchLowercase]}
    {\setmonofont{Latin Modern Mono}[Scale=MatchLowercase]}
\fi
\usepackage[russian]{babel}
\usepackage{geometry}
\geometry{left=25mm,right=15mm,top=20mm,bottom=20mm}
\usepackage{setspace}
\onehalfspacing
\usepackage{indentfirst}
\setlength{\parindent}{1.25cm}
\setlength{\parskip}{2pt}
\setlength{\emergencystretch}{3em}
\usepackage{graphicx,booktabs,longtable,array,tabularx,calc,etoolbox}
\newcounter{none} % Pandoc 3.9 uses this counter for unnumbered longtables.
\usepackage{enumitem}
\setlist[itemize]{leftmargin=1.25cm,itemsep=2pt,topsep=4pt}
\setlist[enumerate]{leftmargin=1.25cm,itemsep=2pt,topsep=4pt}
\newcommand{\tightlist}{\setlength{\itemsep}{0pt}\setlength{\parskip}{0pt}}
\AtBeginEnvironment{longtable}{\small\singlespacing}
\setlength{\LTpre}{8pt}
\setlength{\LTpost}{8pt}
\setlength{\tabcolsep}{5pt}
\setlength{\LTleft}{0pt}
\setlength{\LTright}{0pt}
\usepackage{xurl}
\usepackage[unicode,colorlinks=true,linkcolor=black,urlcolor=blue]{hyperref}
\usepackage{bookmark}
\hypersetup{pdftitle={Вклад Самбьюи в матрицу трассировки},pdfauthor={Самбьюи А.Ж.},pdfsubject={Криптотренажёр mephi-12}}
\urlstyle{same}
\pagestyle{plain}
\setcounter{secnumdepth}{3}
\setcounter{tocdepth}{2}
\listfiles

\begin{document}
\hypersetup{pageanchor=false}
\begin{titlepage}
  \thispagestyle{empty}
  \centering
  {\large НИЯУ МИФИ\par}
  \vspace*{3.5cm}
  {\LARGE\bfseries Вклад Самбьюи в матрицу трассировки\par}
  \vspace{0.5cm}
  {\large Вклад DOC-09\par}
  \vspace{0.6cm}
  {\large Платформа для изучения и практики курса по криптографии\par}
  {\large Криптотренажёр mephi-12\par}
  \vspace{0.7cm}
  {\large Версия 0.1\par}
  {\large Редакция от 25.11.2026\par}
  \vspace{1cm}
  \begin{tabularx}{\textwidth}{@{}p{3.4cm}Xp{4cm}@{}}
    \toprule
    Роль & ФИО & Подпись / Дата \\
    \midrule
    Разработал & Самбьюи А.Ж. & \\
    \addlinespace[6pt]
    Согласовано & & \\
    \addlinespace[6pt]
    Утверждено & & \\
    \bottomrule
  \end{tabularx}
  \vfill
  {\large Москва, 2026\par}
\end{titlepage}
\hypersetup{pageanchor=true}

\section*{Лист регистрации изменений}
\addcontentsline{toc}{section}{Лист регистрации изменений}
\begingroup\small\singlespacing
\noindent
\begin{tabularx}{\textwidth}{@{}p{1.2cm}p{2cm}p{2.5cm}Xp{1.8cm}@{}}
\toprule
Версия & Дата & Автор & Содержание изменения & Утвердил \\
\midrule
0.1 & 25.11.2026 & Самбьюи А.Ж. & Подготовлены связи 13 собственных SR с реальными SWR/TC; общая матрица остаётся за редактором. & \\
\bottomrule
\end{tabularx}
\endgroup
\clearpage
\begingroup\small\singlespacing
\tableofcontents
\endgroup
\clearpage

\hypertarget{doc-09-вклад-самбьюи-в-матрицу-трассировки}{}

Редакция 0.1 от 25.11.2026. Автор --- Самбьюи А.Ж. Статус --- подготовленный доменный вклад для редактора общей матрицы Портнова Л.М. Общий DOC-09 v1.0 этим материалом не объявляется завершённым.

\section{Область и формат}\label{ux43eux431ux43bux430ux441ux442ux44c-ux438-ux444ux43eux440ux43cux430ux442}

\hypertarget{область-и-формат}{}

Включены 13 системных требований: SR-CFG-001--008 и SR-SYS-011--015. Для каждого указаны программное требование и существующий описательный кейс. В 24 строках отражены основные и дополнительные междоменные связи: административный обзор и преподавательская ведомость, назначение группы, повторы, закрепление вариантов и восстановление принятого ответа. Полная редакторская матрица также учитывает остальные общие связи DOC-10. Наличие связи не означает прохождение испытания или принятие требования.

\href{https://github.com/Kokorev-VD/mephi-samolety/blob/main/docs/05_\%D0\%A2\%D1\%80\%D0\%B0\%D1\%81\%D1\%81\%D0\%B8\%D1\%80\%D0\%BE\%D0\%B2\%D0\%BA\%D0\%B0/DOC-09_\%D0\%92\%D0\%BA\%D0\%BB\%D0\%B0\%D0\%B4_\%D0\%A1\%D0\%B0\%D0\%BC\%D0\%B1\%D1\%8C\%D1\%8E\%D0\%B8.csv}{CSV-вклад} использует разделитель «запятая» и UTF-8 с BOM. Поля \texttt{sr\_id}, \texttt{swr\_id}, \texttt{tc\_id} содержат постоянные ID; \texttt{owner} --- владельца вклада; три поля \texttt{*\_ref} --- относительные пути от папки CSV с якорями; \texttt{status} --- фактический статус подготовки; \texttt{relation\_\allowbreak{}note} --- проверяемая часть обязанности. Одинаковый SR может иметь несколько строк и кейсов. Пустых или будущих ID в этой редакции нет.

\section{Связи своего объёма}\label{ux441ux432ux44fux437ux438-ux441ux432ux43eux435ux433ux43e-ux43eux431ux44aux451ux43cux430}

\hypertarget{связи-своего-объёма}{}

{\def\LTcaptype{none} % do not increment counter
\begin{longtable}[]{@{}
  >{\raggedright\arraybackslash}p{(\linewidth - 4\tabcolsep) * \real{0.3333}}
  >{\raggedright\arraybackslash}p{(\linewidth - 4\tabcolsep) * \real{0.3333}}
  >{\raggedright\arraybackslash}p{(\linewidth - 4\tabcolsep) * \real{0.3333}}@{}}
\toprule\noalign{}
\begin{minipage}[b]{\linewidth}\raggedright
Системное требование
\end{minipage} & \begin{minipage}[b]{\linewidth}\raggedright
Программное требование
\end{minipage} & \begin{minipage}[b]{\linewidth}\raggedright
Описательный кейс
\end{minipage} \\
\midrule\noalign{}
\endhead
\bottomrule\noalign{}
\endlastfoot
{\href{../04\_Системные\_требования/DOC-05\_Системные\_требования.pdf\#sr-cfg-001}{SR-CFG-001}} & {\href{../07\_Программные\_требования/DOC-06.6\_Конфигурации\_и\_модерация.pdf\#swr-cfg-001}{SWR-CFG-001}} & {\href{../10\_Верификация\_и\_тестирование/DOC-10\_План\_верификации.pdf\#tc-cfg-001}{TC-CFG-001}} \\
{\href{../04\_Системные\_требования/DOC-05\_Системные\_требования.pdf\#sr-cfg-002}{SR-CFG-002}} & {\href{../07\_Программные\_требования/DOC-06.6\_Конфигурации\_и\_модерация.pdf\#swr-cfg-002}{SWR-CFG-002}} & {\href{../10\_Верификация\_и\_тестирование/DOC-10\_План\_верификации.pdf\#tc-cfg-002}{TC-CFG-002}} \\
{\href{../04\_Системные\_требования/DOC-05\_Системные\_требования.pdf\#sr-cfg-003}{SR-CFG-003}} & {\href{../07\_Программные\_требования/DOC-06.6\_Конфигурации\_и\_модерация.pdf\#swr-cfg-003}{SWR-CFG-003}} & {\href{../10\_Верификация\_и\_тестирование/DOC-10\_План\_верификации.pdf\#tc-cfg-003}{TC-CFG-003}} \\
{\href{../04\_Системные\_требования/DOC-05\_Системные\_требования.pdf\#sr-cfg-004}{SR-CFG-004}} & {\href{../07\_Программные\_требования/DOC-06.6\_Конфигурации\_и\_модерация.pdf\#swr-cfg-004}{SWR-CFG-004}} & {\href{../10\_Верификация\_и\_тестирование/DOC-10\_План\_верификации.pdf\#tc-cfg-004}{TC-CFG-004}} \\
{\href{../04\_Системные\_требования/DOC-05\_Системные\_требования.pdf\#sr-cfg-005}{SR-CFG-005}} & {\href{../07\_Программные\_требования/DOC-06.6\_Конфигурации\_и\_модерация.pdf\#swr-cfg-005}{SWR-CFG-005}} & {\href{../10\_Верификация\_и\_тестирование/DOC-10\_План\_верификации.pdf\#tc-cfg-005}{TC-CFG-005}} \\
{\href{../04\_Системные\_требования/DOC-05\_Системные\_требования.pdf\#sr-cfg-006}{SR-CFG-006}} & {\href{../07\_Программные\_требования/DOC-06.6\_Конфигурации\_и\_модерация.pdf\#swr-cfg-006}{SWR-CFG-006}} & {\href{../10\_Верификация\_и\_тестирование/DOC-10\_План\_верификации.pdf\#tc-cfg-006}{TC-CFG-006}} \\
{\href{../04\_Системные\_требования/DOC-05\_Системные\_требования.pdf\#sr-cfg-007}{SR-CFG-007}} & {\href{../07\_Программные\_требования/DOC-06.6\_Конфигурации\_и\_модерация.pdf\#swr-cfg-007}{SWR-CFG-007}} & {\href{../10\_Верификация\_и\_тестирование/DOC-10\_План\_верификации.pdf\#tc-cfg-007}{TC-CFG-007}} \\
{\href{../04\_Системные\_требования/DOC-05\_Системные\_требования.pdf\#sr-cfg-008}{SR-CFG-008}} & {\href{../07\_Программные\_требования/DOC-06.6\_Конфигурации\_и\_модерация.pdf\#swr-cfg-008}{SWR-CFG-008}} & {\href{../10\_Верификация\_и\_тестирование/DOC-10\_План\_верификации.pdf\#tc-cfg-008}{TC-CFG-008}} \\
{\href{../04\_Системные\_требования/DOC-05\_Системные\_требования.pdf\#sr-sys-011}{SR-SYS-011}} & {\href{../07\_Программные\_требования/DOC-06\_Программные\_требования.pdf\#swr-sys-010}{SWR-SYS-010}} & {\href{../10\_Верификация\_и\_тестирование/DOC-10\_План\_верификации.pdf\#tc-sys-002}{TC-SYS-002}} \\
{\href{../04\_Системные\_требования/DOC-05\_Системные\_требования.pdf\#sr-sys-012}{SR-SYS-012}} & {\href{../07\_Программные\_требования/DOC-06\_Программные\_требования.pdf\#swr-sys-011}{SWR-SYS-011}} & {\href{../10\_Верификация\_и\_тестирование/DOC-10\_План\_верификации.pdf\#tc-sys-003}{TC-SYS-003}} \\
{\href{../04\_Системные\_требования/DOC-05\_Системные\_требования.pdf\#sr-sys-013}{SR-SYS-013}} & {\href{../07\_Программные\_требования/DOC-06\_Программные\_требования.pdf\#swr-sys-012}{SWR-SYS-012}} & {\href{../10\_Верификация\_и\_тестирование/DOC-10\_План\_верификации.pdf\#tc-sys-003}{TC-SYS-003}} \\
{\href{../04\_Системные\_требования/DOC-05\_Системные\_требования.pdf\#sr-sys-014}{SR-SYS-014}} & {\href{../07\_Программные\_требования/DOC-06\_Программные\_требования.pdf\#swr-sys-013}{SWR-SYS-013}} & {\href{../10\_Верификация\_и\_тестирование/DOC-10\_План\_верификации.pdf\#tc-sys-004}{TC-SYS-004}} \\
{\href{../04\_Системные\_требования/DOC-05\_Системные\_требования.pdf\#sr-sys-015}{SR-SYS-015}} & {\href{../07\_Программные\_требования/DOC-06\_Программные\_требования.pdf\#swr-sys-014}{SWR-SYS-014}} & {\href{../10\_Верификация\_и\_тестирование/DOC-10\_План\_верификации.pdf\#tc-sys-005}{TC-SYS-005}} \\
{\href{../04\_Системные\_требования/DOC-05\_Системные\_требования.pdf\#sr-cfg-001}{SR-CFG-001}} & {\href{../07\_Программные\_требования/DOC-06.2\_Сессии\_задач\_и\_планировщик.pdf\#swr-ses-002}{SWR-SES-002}} & {\href{../10\_Верификация\_и\_тестирование/DOC-10\_План\_верификации.pdf\#tc-cfg-001}{TC-CFG-001}} \\
{\href{../04\_Системные\_требования/DOC-05\_Системные\_требования.pdf\#sr-cfg-002}{SR-CFG-002}} & {\href{../07\_Программные\_требования/DOC-06\_Программные\_требования.pdf\#swr-auth-009}{SWR-AUTH-009}} & {\href{../10\_Верификация\_и\_тестирование/DOC-10\_План\_верификации.pdf\#tc-auth-005}{TC-AUTH-005}} \\
{\href{../04\_Системные\_требования/DOC-05\_Системные\_требования.pdf\#sr-cfg-004}{SR-CFG-004}} & {\href{../07\_Программные\_требования/DOC-06\_Программные\_требования.pdf\#swr-auth-009}{SWR-AUTH-009}} & {\href{../10\_Верификация\_и\_тестирование/DOC-10\_План\_верификации.pdf\#tc-auth-005}{TC-AUTH-005}} \\
{\href{../04\_Системные\_требования/DOC-05\_Системные\_требования.pdf\#sr-cfg-005}{SR-CFG-005}} & {\href{../07\_Программные\_требования/DOC-06.2\_Сессии\_задач\_и\_планировщик.pdf\#swr-ses-003}{SWR-SES-003}} & {\href{../10\_Верификация\_и\_тестирование/DOC-10\_План\_верификации.pdf\#tc-ses-003}{TC-SES-003}} \\
{\href{../04\_Системные\_требования/DOC-05\_Системные\_требования.pdf\#sr-cfg-005}{SR-CFG-005}} & {\href{../07\_Программные\_требования/DOC-06.2\_Сессии\_задач\_и\_планировщик.pdf\#swr-ses-002}{SWR-SES-002}} & {\href{../10\_Верификация\_и\_тестирование/DOC-10\_План\_верификации.pdf\#tc-cfg-005}{TC-CFG-005}} \\
{\href{../04\_Системные\_требования/DOC-05\_Системные\_требования.pdf\#sr-cfg-007}{SR-CFG-007}} & {\href{../07\_Программные\_требования/DOC-06.2\_Сессии\_задач\_и\_планировщик.pdf\#swr-ses-019}{SWR-SES-019}} & {\href{../10\_Верификация\_и\_тестирование/DOC-10\_План\_верификации.pdf\#tc-cfg-007}{TC-CFG-007}} \\
{\href{../04\_Системные\_требования/DOC-05\_Системные\_требования.pdf\#sr-cfg-008}{SR-CFG-008}} & {\href{../07\_Программные\_требования/DOC-06.6\_Конфигурации\_и\_модерация.pdf\#swr-cfg-005}{SWR-CFG-005}} & {\href{../10\_Верификация\_и\_тестирование/DOC-10\_План\_верификации.pdf\#tc-cfg-005}{TC-CFG-005}} \\
{\href{../04\_Системные\_требования/DOC-05\_Системные\_требования.pdf\#sr-cfg-008}{SR-CFG-008}} & {\href{../07\_Программные\_требования/DOC-06.2\_Сессии\_задач\_и\_планировщик.pdf\#swr-ses-003}{SWR-SES-003}} & {\href{../10\_Верификация\_и\_тестирование/DOC-10\_План\_верификации.pdf\#tc-ses-003}{TC-SES-003}} \\
{\href{../04\_Системные\_требования/DOC-05\_Системные\_требования.pdf\#sr-cfg-008}{SR-CFG-008}} & {\href{../07\_Программные\_требования/DOC-06.6\_Конфигурации\_и\_модерация.pdf\#swr-cfg-007}{SWR-CFG-007}} & {\href{../10\_Верификация\_и\_тестирование/DOC-10\_План\_верификации.pdf\#tc-cfg-007}{TC-CFG-007}} \\
{\href{../04\_Системные\_требования/DOC-05\_Системные\_требования.pdf\#sr-sys-012}{SR-SYS-012}} & {\href{../07\_Программные\_требования/DOC-06\_Программные\_требования.pdf\#swr-sys-008}{SWR-SYS-008}} & {\href{../10\_Верификация\_и\_тестирование/DOC-10\_План\_верификации.pdf\#tc-sys-003}{TC-SYS-003}} \\
{\href{../04\_Системные\_требования/DOC-05\_Системные\_требования.pdf\#sr-sys-012}{SR-SYS-012}} & {\href{../07\_Программные\_требования/DOC-06.2\_Сессии\_задач\_и\_планировщик.pdf\#swr-ses-021}{SWR-SES-021}} & {\href{../10\_Верификация\_и\_тестирование/DOC-10\_План\_верификации.pdf\#tc-ses-013}{TC-SES-013}} \\
\end{longtable}
}

\section{Передача редактору и контроль}\label{ux43fux435ux440ux435ux434ux430ux447ux430-ux440ux435ux434ux430ux43aux442ux43eux440ux443-ux438-ux43aux43eux43dux442ux440ux43eux43bux44c}

\hypertarget{передача-редактору-и-контроль}{}

При включении в общую матрицу редактор сохраняет ID, объединяет дополнительные связи из DOC-10 и пересчитывает покрытие по полному перечню SR всех владельцев. Для собственного вклада сверяют существование определений, смысл соответствия и наличие условий, действий и ожидаемого результата в кейсе. TC-SYS-003 покрывает и сохранность, и готовность запуска; совпадение TC в двух строках допустимо. 13 из 13 означает наличие связей для 13 требований этого вклада, а не доказательство достаточности одним счётом строк. Преподавательская часть SR-CFG-002 прослежена отдельно через SWR-AUTH-009 и TC-AUTH-005; повторная выдача SR-CFG-008 --- через SWR-CFG-005, SWR-SES-003 и соответствующие кейсы. Смысловое покрытие требует прочитать все связанные обязанности и случаи.

Предварительная сверка источников выполнена при подготовке. Формальное принятие общей матрицы и результаты исполнения не подтверждены. Закрытие замечаний ФИ-СТ зависит от появления реального DOC-13.3; его отсутствие не заменяется отметкой в CSV.


\end{document}
